%% sec6_stats.tex -- Section 6: Statistical Analysis

\section{Statistical Analysis}
\label{sec:stats}

The paired experimental design yields, for each value of $n$ and each trial
$b \in [B]$, a triple $(\hat{p}^{\mathrm{rand}}_b, \hat{p}^{\mathrm{gr}}_b,
\hat{p}^{\mathrm{SA}}_b)$ evaluated against the same $50{,}000$ draws.
The within-trial differences $d_b = \hat{p}^{\mathrm{SA}}_b -
\hat{p}^{\mathrm{rand}}_b$ are i.i.d.\ across trials and are tested directly
by a paired one-sample $t$-test of $H_0: E[d_b] = 0$ against $H_1: E[d_b] > 0$.
Cohen's $d$ effect size \cite{Cohen1988} is defined as
$d_{\mathrm{eff}} = \bar{d}/s_d$, where $\bar{d}$ and $s_d$ are the
sample mean and standard deviation of $\{d_b\}$.

Table~\ref{tab:ttest} reports the results.  All values $n \ge 2$ yield
$p < 10^{-8}$, and Cohen's $d$ ranges from $1.18$ at $n=2$ to $3.50$ at
$n=7$.  Any threshold of practical significance (small: 0.2, medium: 0.5,
large: 0.8) is exceeded by a wide margin at every tested value of $n$.

\begin{table}[ht]
  \centering
  \caption{Paired $t$-test of Greedy+SA vs.\ Random.
           $\bar{d}$: mean gain (pp).
           $t_n$: test statistic.
           $p$: two-tailed $p$-value.
           $d_{\mathrm{eff}}$: Cohen's $d$ \cite{Cohen1988}.}
  \label{tab:ttest}
  \smallskip
  \begin{tabular}{@{}rrrrrrl@{}}
    \toprule
    $n$ & $B$ & $\bar{d}$ (pp) & $t_n$ & $p$ & $d_{\mathrm{eff}}$ & sig \\
    \midrule
    1  & 50 & $-0.03$ &  $-1.33$ & $1.9 \times 10^{-1}$ & $-0.19$ & n.s. \\
    2  & 50 & $+6.2$  &   $8.31$ & $6.4 \times 10^{-11}$ & $1.18$ & *** \\
    3  & 50 & $+8.1$  &  $13.26$ & $7.9 \times 10^{-18}$ & $1.88$ & *** \\
    5  & 50 & $+11.0$ &  $18.66$ & $6.6 \times 10^{-24}$ & $2.64$ & *** \\
    7  & 40 & $+10.3$ &  $22.16$ & $1.1 \times 10^{-23}$ & $3.50$ & *** \\
    10 & 40 & $+8.3$  &  $15.82$ & $1.5 \times 10^{-18}$ & $2.50$ & *** \\
    15 & 30 & $+7.0$  &  $14.55$ & $7.3 \times 10^{-15}$ & $2.66$ & *** \\
    20 & 30 & $+4.4$  &  $17.68$ & $4.5 \times 10^{-17}$ & $3.23$ & *** \\
    25 & 30 & $+2.9$  &  $13.12$ & $1.0 \times 10^{-13}$ & $2.40$ & *** \\
    30 & 30 & $+2.0$  &  $17.06$ & $1.2 \times 10^{-16}$ & $3.12$ & *** \\
    40 & 30 & $+0.7$  &  $10.40$ & $2.7 \times 10^{-11}$ & $1.90$ & *** \\
    50 & 30 & $+0.2$  &   $8.32$ & $3.6 \times 10^{-9}$  & $1.52$ & *** \\
    75 & 30 & $+0.02$ &   $7.18$ & $6.7 \times 10^{-8}$  & $1.31$ & *** \\
    100& 30 & $+0.00$ &   $4.01$ & $3.9 \times 10^{-4}$  & $0.73$ & *** \\
    \bottomrule
    \multicolumn{7}{@{}l@{}}{\footnotesize
      *** $p < 0.001$; n.s.\ not significant.}
  \end{tabular}
\end{table}

The absolute gain $\bar{d}$ peaks at $n=5$ ($+11.0$ pp) and remains above
ten percentage points at $n=7$.  The effect size $d_{\mathrm{eff}}$ peaks at
$n=7$ ($3.50$) and stays well above conventional large-effect thresholds
throughout $n \le 75$.  A notable asymmetry is that the standard deviation
of $\hat{p}^{\mathrm{rand}}$ is three to four times larger than that of
$\hat{p}^{\mathrm{SA}}$: the optimised strategy produces higher expected
probability and more predictable outcomes.

Defining the signal-to-noise ratio as
$\mathrm{SNR}(n) = t_n / \sqrt{B}$, we find values ranging from
$8\times$ at $n=2$ to $22\times$ at $n=7$.  An SNR above $3\times$
(roughly corresponding to $p < 0.003$ in a normal test) is maintained
throughout $n \le 100$, confirming that the reported gains would be
detectable with substantially fewer trials.

For a practical saturation characterisation, define $n^*_\epsilon$ as the
smallest $n$ for which $P_{\mathrm{win}}^{\mathrm{SA}}(n) \ge 1-\epsilon$.
From Table~\ref{tab:pwin-full}, $n^*_{0.001} \approx 75$ and
$n^*_{0.01} \approx 50$ for both strategies.  The meaningful range for
optimisation is $n \in [2, 40]$, where the gain exceeds $0.7$ pp and
$d_{\mathrm{eff}} > 1.9$.  Within this range, $n \in [3, 10]$ offers the
best return: gains above eight percentage points with very large effect sizes.
